% %%%%%%%%%%%%%%%%%%%%%%%%%%%%%%%%%%%%%%%%%%%%%%%%%%%%%%%%%%%%%%%%%%%%%
% %% State-of-the-Art %%
% Conclusion
% %%%%%%%%%%%%%%%%%%%%%%%%%%%%%%%%%%%%%%%%%%%%%%%%%%%%%%%%%%%%%%%%%%%%%

\section{Conclusion}
\label{sec:soa:conclu}

\lettrine{T}{he} rapid expansion of data-intensive applications in computational physics, machine learning, and graph analytics has driven the development of numerous hardware accelerators over the past ten years, each aiming to surpass the performance limitations of traditional CPUs and GPUs.
This chapter presented a comprehensive review of the most significant designs, offering a structured comparison and highlighting the core hardware challenges associated with sparse matrix multiplication kernels.

Our analysis reveals that the primary challenge in sparse matrix computations lies in efficiently managing memory access patterns, particularly the dereferencing required for \glspl{spFormat}.
Depending on the chosen algorithm, \acrfull{ip}, \acrfull{op}, or \acrlong{gust}'s (\acrshort{gust}), the main bottleneck shifts between \gls{gather} and \gls{intersection} operations (for \acrshort{ip}) or \gls{scatter} and \gls{reduction} operations (for \acrshort{op}).
\acrlong{gust}'s algorithm emerges as a balanced solution, as it reduces the number of $B$-rows that need to be accessed while generating outputs row-by-row, simplifying both \gls{gather} and \gls{scatter} processes.
The most advanced \acrlong{gust}-based accelerators further optimize performance through intelligent cache \gls{eviction} strategies, where victims are selected based on forward-looking analysis of $A$-indices.

While \textit{standalone} accelerators dominate the landscape, \textit{processor-assist} designs, such as those that modify cache behavior or leverage existing memory resources, also play a crucial role, particularly in systems where flexibility and lower area overhead are prioritized.
In contrast, high-performance \textit{standalone} accelerators often feature dedicated, massively parallel \gls{reduction}/merger units to maximize parallelism and minimize memory bottlenecks.

The diversity of existing designs underscores that there is no single ``best'' sparse accelerator.
Instead, the optimal architecture depends on the specific problem class, matrix structure, and performance requirements.
Future advancements will likely focus on tailoring accelerators to the unique characteristics of different matrix structures, ensuring that caches, memory hierarchies, and computational units are optimized for the workload at hand.

This chapter not only synthesizes the current state of sparse matrix acceleration but also provides a roadmap for future research, emphasizing the need for architectures that are both highly efficient and adaptable to the diverse and evolving landscape of sparse data processing.

\bigskip\bigskip

\begin{highlight}{Takeaways}{Mulberry}
    \colorbf{Mulberry}{Classification of Sparse Matrix Accelerators:}
    Hardware accelerators for sparse matrix computations can be categorized by:
    \begin{itemize}
        \item \textbf{Autonomy:} \textit{standalone} versus \textit{processor-assist}
        \item \textbf{Position in the memory hierarchy:} near-processor, near-cache, or near-main-memory
        \item \textbf{Algorithm compatibility:} \acrfull{ip}, \acrfull{op}, or \acrfull{gust}
        \item \textbf{Supported Memory Flows:} \gls{gather}/\gls{scatter} hardware support
        \item \textbf{Targeted kernels:} sparse \acrfull{mv}, sparse \acrfull{mm}, or both
    \end{itemize}

    \bigskip
    \colorbf{Mulberry}{Key Findings:}
    \begin{itemize}
        \item \textbf{Algorithm Selection:} \acrlong{gust}'s algorithm offers the best trade-off for hardware implementation, balancing input reuse and output generation, while the \acrlong{op} algorithm remains a strong alternative, particularly for parallelism.
        \item \textbf{Memory Optimization:} Efficient use of local memory, whether through custom caches or smart \gls{eviction} policies, is critical to reduce off-chip traffic and to improve performance.
        \item \textbf{Adaptability:} The most successful accelerators dynamically adjust their dataflows and memory strategies to match the sparsity patterns of the input matrices, highlighting the importance of adaptability.
        \item \textbf{Scalability and Parallelism:} High-performance designs leverage parallel \gls{reduction} and merger units and optimized memory hierarchies.
    \end{itemize}

    \bigskip
    \colorbf{Mulberry}{Future Research Directions:}
    Next-generation accelerators will increasingly feature \textbf{adaptive mechanisms}, \textbf{multi-core architectures}, and \textbf{specialized memory systems} tailored to diverse sparse matrix structures and evolving application requirements.

    \bigskip
    \colorbf{Mulberry}{SpDCache:}
    Our work focuses on near-cache, \textit{processor-assist} accelerators that compute \gls{scatter}-focused \acrlong{op} \acrshort{spmv} kernels.
\end{highlight}

\clearpage